\chapter{Setup: hardware, Pi, and an empty project}
\label{ch:setup}

\chapterbrief
  {An RC522 wired to a Pi, a Java project that compiles, and one green test.}
  {None yet --- this chapter removes every excuse the environment can give you later.}
  {\code{mvn compile} passing, and a chip you can talk to.}

Nothing in this chapter is interesting, and that is the point. Every hour spent
later wondering \emph{``is this my code or my wiring?''} is an hour you can buy
back now by making the environment boring and trustworthy.

\section{Wiring the RC522}
\index{RC522!wiring}\index{SPI!wiring}

The RC522 speaks SPI: six signal wires plus power.

\begin{center}\small
\begin{tabular}{@{}llll@{}}
\toprule
\thead{RC522 pin} & \thead{Pi pin} & \thead{Pi name} & \thead{Why} \\
\midrule
\reg{3.3V} & 1  & 3V3            & \textbf{Never 5\,V.} The MFRC522 is a 3.3\,V part. \\
\reg{RST}  & 22 & BCM 25         & Hardware reset \\
\reg{GND}  & 6  & GND            & \\
\reg{MISO} & 21 & BCM 9 (MISO)   & Chip $\rightarrow$ Pi \\
\reg{MOSI} & 19 & BCM 10 (MOSI)  & Pi $\rightarrow$ chip \\
\reg{SCK}  & 23 & BCM 11 (SCLK)  & Clock \\
\reg{SDA}  & 24 & BCM 8 (CE0)    & Chip select --- confusingly labelled with an I\textsuperscript{2}C name \\
\reg{IRQ}  & --- & not connected & We poll instead; simpler, and fast enough \\
\bottomrule
\end{tabular}
\end{center}

\begin{center}
\begin{tikzpicture}[
   pin/.style={font=\scriptsize\sffamily},
   wire/.style={line width=0.9pt},
   x=1cm,y=1cm]

  % Pi header
  \draw[rounded corners=2pt, brewernavy, line width=0.8pt] (0,0) rectangle (1.5,5.2);
  \node[pin, color=brewernavy, rotate=90] at (0.32,2.6) {Raspberry Pi header};
  \foreach \y/\lbl in {4.7/1, 4.1/6, 3.5/19, 2.9/21, 2.3/22, 1.7/23, 1.1/24}{
     \fill[brewernavy!70] (1.05,\y) circle (1.6pt);
     \node[pin, color=brewergrey, anchor=east] at (1.0,\y) {\lbl};
  }

  % RC522
  \draw[rounded corners=2pt, breweramber!80!black, line width=0.8pt] (6.0,0.4) rectangle (8.4,5.0);
  \node[pin, color=breweramber!70!black, rotate=90] at (8.1,2.7) {RC522};
  \foreach \y/\lbl in {4.7/3.3V, 4.1/GND, 3.5/MOSI, 2.9/MISO, 2.3/RST, 1.7/SCK, 1.1/SDA}{
     \fill[breweramber!80!black] (6.15,\y) circle (1.6pt);
     \node[pin, color=brewergrey, anchor=west] at (6.3,\y) {\lbl};
  }

  % wires
  \foreach \y/\c in {4.7/brewerred, 4.1/brewerink, 3.5/brewerblue,
                     2.9/brewergreen, 2.3/breweramber, 1.7/brewerblue, 1.1/brewergreen}{
     \draw[wire, color=\c!75] (1.05,\y) -- (6.15,\y);
  }
  \node[pin, color=brewergrey] at (3.6,5.35) {short wires --- this matters};
\end{tikzpicture}
\end{center}

\begin{trap}{Long wires are not a style question}
At 1\,MHz over 20\,cm of loose DuPont jumper you get intermittent failures that
look exactly like firmware bugs: a version register that reads \hex{0x00} on one
run and \hex{0x92} on the next.

If the chip is flaky, \textbf{shorten the wires before you debug a single line
of code}. More people have lost an evening to this than to any bug in this book.
\end{trap}

\section{Pi OS and SPI}

Flash Raspberry Pi OS (64-bit), then enable SPI:

\begin{shell}
sudo raspi-config       # Interface Options -> SPI -> Enable
sudo reboot
ls -l /dev/spidev0.0    # must exist after the reboot
\end{shell}

\code{/dev/spidev0.0} is the kernel exposing ``SPI bus 0, chip-select 0'' as a
file. That file is the entire hardware interface: everything in Part~II is
ultimately reads and writes to it.

Add your user to the hardware groups, or Pi4J cannot open the device:

\begin{shell}
sudo usermod -aG spi,gpio $USER    # log out and back in
sudo apt update
sudo apt install -y openjdk-21-jdk maven
java -version                      # expect 21.x
\end{shell}

\section{Create the project}

Develop on a laptop and deploy to the Pi. Compiling on a Zero 2 W is slow
enough to break your concentration.

\begin{shell}
mkdir midnight-brewer && cd midnight-brewer
git init
mkdir -p src/main/java/com/midnightbrewer/hardware
mkdir -p src/test/java/com/midnightbrewer/hardware
\end{shell}

\subsection{The \texttt{pom.xml}}

Maven's project file. Three things matter: Java 21, JUnit 5, and Pi4J.

\begin{xcode}
<project xmlns="http://maven.apache.org/POM/4.0.0">
  <modelVersion>4.0.0</modelVersion>
  <groupId>com.midnightbrewer</groupId>
  <artifactId>MidnightBrewer</artifactId>
  <version>1.0-SNAPSHOT</version>

  <properties>
    <maven.compiler.source>21</maven.compiler.source>
    <maven.compiler.target>21</maven.compiler.target>
    <project.build.sourceEncoding>UTF-8</project.build.sourceEncoding>
    <pi4j.version>3.0.2</pi4j.version>
  </properties>

  <dependencies>
    <!-- Pi4J: Java access to the Pi's SPI and GPIO -->
    <dependency>
      <groupId>com.pi4j</groupId><artifactId>pi4j-core</artifactId>
      <version>${pi4j.version}</version>
    </dependency>
    <!-- the SPI provider: talks to /dev/spidev0.0 -->
    <dependency>
      <groupId>com.pi4j</groupId><artifactId>pi4j-plugin-linuxfs</artifactId>
      <version>${pi4j.version}</version>
    </dependency>
    <!-- the GPIO provider, for the reset pin -->
    <dependency>
      <groupId>com.pi4j</groupId><artifactId>pi4j-plugin-gpiod</artifactId>
      <version>${pi4j.version}</version>
    </dependency>
    <!-- Pi4J logs through SLF4J; without a binding it prints warnings -->
    <dependency>
      <groupId>org.slf4j</groupId><artifactId>slf4j-simple</artifactId>
      <version>2.0.13</version>
    </dependency>

    <dependency>
      <groupId>org.junit.jupiter</groupId><artifactId>junit-jupiter</artifactId>
      <version>5.10.2</version>
      <scope>test</scope>
    </dependency>
  </dependencies>

  <build><plugins>
    <plugin>
      <groupId>org.apache.maven.plugins</groupId>
      <artifactId>maven-surefire-plugin</artifactId>
      <version>3.2.5</version>
    </plugin>
  </plugins></build>
</project>
\end{xcode}

\begin{concept}{\texttt{<scope>test</scope>} is a boundary, not a detail}
It means ``this library exists when compiling and running tests, but is not
shipped with the application.'' Get it wrong and your kiosk drags a test
framework onto the Pi.

More usefully, it \emph{enforces} the boundary: if you accidentally
\code{import org.junit} in main code, the build fails. A dependency scope is the
build system letting you state a design rule and have it checked.
\end{concept}

\begin{shell}
mvn -q compile     # should print nothing and exit 0
\end{shell}

There is nothing to compile yet. You have proven that Maven, Java and the
network can all do their jobs --- which is exactly one fewer thing that can be
wrong tomorrow.

\section{Java for a C programmer}
\index{Java!versus C}

You will meet the rest as you go, but four things trip up every C programmer on
the first day.

\subsection{Every function lives inside a class}

There are no free functions. \code{Math.abs()} is a function \code{abs} inside a
class \code{Math}. A file \code{Foo.java} contains \code{class Foo}. This feels
bureaucratic, and for small programs it is --- but it is also why the language
can make guarantees C cannot, and Part~II is where those guarantees start paying
for themselves.

\subsection{\texttt{byte} is signed, and there is no \texttt{unsigned}}
\index{signed byte}

This will bite you within the hour.

\begin{jcode}
byte b = (byte) 0x92;
System.out.println(b);          // -110   <-- the same bits, read as signed
System.out.println(b & 0xFF);   // 146    <-- what you meant
\end{jcode}

\hex{0x92} is \code{1001 0010}. With the top bit set, Java reads it as negative.
The fix is the idiom \code{\& 0xFF}, which promotes to \code{int} and masks off
the sign extension.

\begin{trap}{The \texttt{\& 0xFF} habit}
\textbf{Any time you turn a \code{byte} into a number, write \code{\& 0xFF}.}
You will forget this at least once and lose an hour to a register that reads
$-110$ instead of \hex{0x92}. Everyone does. Now you have been warned, and you
will still do it.
\end{trap}

\subsection{Arrays know their own length, and are bounds-checked}

\code{arr.length} --- no parentheses, it is a field. Reading past the end throws
an exception naming the line, rather than silently returning whatever was next in
memory. This converts C's worst class of bug into a stack trace.

\subsection{No manual memory management}

No \code{malloc}, no \code{free}. \code{new byte[16]} gives you sixteen
\emph{zeroed} bytes, and they vanish when nothing refers to them. The zeroing
matters more than you would think: several structures in this book are ``a
16-byte buffer that is mostly zeros'', and you get that for free.

\medskip
One C habit transfers perfectly: bit manipulation. \code{<<}, \code{>>},
\code{\&}, \code{|}, \code{\^{}} and \code{\~{}} all mean exactly what they mean
in C. The register-poking half of this project will feel like home.

\section{Commit}

\begin{shell}
git add -A
git commit -m "chore: empty Maven project, Java 21, Pi4J + JUnit"
\end{shell}
